% Teil II: durchgerechneter Weg, Teile a--b.

\subsection*{Gegeben und gesucht}
$a)$ Zwei Geschosse gleicher kinetischer Energie, eines mit doppelter
Masse. $b)$ Körper $m=\qty{3}{kg}$ aus der Ruhe gleichmässig
beschleunigt, nach $\qty{6}{s}$ Geschwindigkeit $\qty{30}{m/s}$.
Gesucht: $a)$ welches schneller (Faktor) und welches mehr Arbeit;
$b)$ Momentanleistung nach $\qty{3}{s}$.

\subsection*{Ansatz}
$a)$ $E_{\mathrm{kin}}=\tfrac12mv^2$ gleichsetzen. $b)$ $a$ aus
$v=at$, dann $P(3\,\mathrm s)=m a v(3\,\mathrm s)$.

\subsection*{Rechnung}
\textbf{a) Geschwindigkeiten.} Sei Masse $m$ (Geschoss $v_1$) und Masse
$2m$ (Geschoss $v_2$); gleiche Energie:
\begin{equation}
  \tfrac12 m v_1^2=\tfrac12 (2m)v_2^2\;\Rightarrow\;
  v_1^2=2v_2^2\;\Rightarrow\;
  \boxed{\;v_1=\sqrt2\,v_2\;}.
\end{equation}
Das leichtere Geschoss ist um den Faktor $\sqrt2\approx1{,}41$ schneller.
Da beide die \emph{gleiche} kinetische Energie und damit dieselbe
Arbeitsfaehigkeit haben, können beide gleich viel Arbeit verrichten.

\textbf{b) Momentanleistung.} Aus $v(t)=at$ mit $v(\qty{6}{s})=\qty{30}{m/s}$:
\begin{equation}
  a=\frac{v}{t_{\mathrm{end}}}=\frac{\qty{30}{m/s}}{\qty{6}{s}}=\qty{5}{m/s^2}.
\end{equation}
Nach $\qty{3}{s}$: $v_3=at_3=\qty{5}{m/s^2}\cdot\qty{3}{s}=\qty{15}{m/s}$. Damit
\begin{equation}
  P=m\,a\,v_3=\qty{3}{kg}\cdot\qty{5}{m/s^2}\cdot\qty{15}{m/s}=
  \boxed{\;\qty{225}{W}\;}.
\end{equation}

\begin{ergebnis}
$a)$ $v_1=\sqrt2\,v_2$ (das leichtere verdoppelt seine Geschwindigkeit um
$\sqrt2$); gleiche kinetische Energie, also gleiche Arbeitsfaehigkeit.
$b)$ $P(\qty{3}{s})=\qty{225}{W}$.
\end{ergebnis}

\begin{kontrolle}
\textbf{Einheiten:} $\mathrm{kg\cdot m/s^2\cdot m/s}=\mathrm{kg\cdot m^2/s^3}
=\unit{\watt}$. \checkmark

\textbf{Querprobe:} Bei $t=\qty{6}{s}$ ist $P=\qty{3}{kg}\cdot\qty{5}{m/s^2}\cdot
\qty{30}{m/s}=\qty{450}{W}$, doppelt so hoch wie bei $\qty{3}{s}$ (weil $P\propto
t$); längere Beschleunigung gibt mehr Leistung. \checkmark
\end{kontrolle}

\begin{merksatz}
$E_{\mathrm{kin}}=\tfrac12mv^2$: die Geschwindigkeit zählt quadratisch.
Momentanleistung $P=mav$; bei gleichmässiger Beschleunigung aus der Ruhe
wächst sie linear mit der Zeit.
\end{merksatz}
